Дискретизация уравнения Рейнольдса (раздел~3.2) при полном заполнении зазора ($\vartheta = 1$, стационарный случай) приводит к линейной системе~\eqref{eq:5point_stencil} с пятиточечным шаблоном. Матрица этой системы разреженная; для внутренних узлов каждая строка содержит не более пяти ненулевых коэффициентов (5-точечный шаблон). Число неизвестных определяется числом внутренних узлов сетки~$N_1 \times N_2$. Коэффициенты $a_E$,~$a_W$,~$a_N$,~$a_S$ неотрицательны, поскольку проводимость $K = h^3/(12\mu) > 0$, а диагональное преобладание $a_P = a_E + a_W + a_N + a_S$ обеспечивает благоприятные свойства системы и способствует сходимости стационарных итерационных методов. Прямые методы (например, LU-разложение) для больших сеток неэффективны из-за заполнения разреженной структуры, поэтому используются итерационные методы, допускающие поточечное обновление.

\textit{Метод Гаусса–Зейделя.}
Метод Якоби (одновременное обновление всех узлов по значениям предыдущей итерации) для уравнения Рейнольдса практически не применяется из-за медленной сходимости; его естественное улучшение~-- метод Гаусса–Зейделя. Формула обновления имеет вид
\begin{equation}\label{eq:gs_update}
  p_{i,j}^{(m+1)} = \frac{1}{a_P}\bigl(
    a_E\,p_{i+1,j}^{(m)} + a_W\,p_{i-1,j}^{(m+1)}
    + a_N\,p_{i,j+1}^{(m)} + a_S\,p_{i,j-1}^{(m+1)}
    + b_{i,j}\bigr),
\end{equation}
\noindent где верхний индекс~$(m)$~-- номер итерации. При последовательном обходе сетки по строкам значения $p_{i-1,j}$ и $p_{i,j-1}$ уже обновлены на текущей итерации (индекс~$m+1$), что ускоряет сходимость по сравнению с методом Якоби. Узлы с условиями Дирихле не обновляются~-- их значения фиксированы (раздел~3.1). Периодичность учитывается связыванием пар узлов на противоположных границах расчётной области.

\textit{Последовательная верхняя релаксация (SOR).}
Для ускорения сходимости вводится параметр релаксации~$\omega$. Обновление выполняется в два шага:
\begin{equation}\label{eq:sor_update}
\begin{split}
  &\tilde{p}_{i,j} = \frac{1}{a_P}\bigl(
    a_E\,p_{i+1,j}^{(m)} + a_W\,p_{i-1,j}^{(m+1)}
    + a_N\,p_{i,j+1}^{(m)} + a_S\,p_{i,j-1}^{(m+1)}
    + b_{i,j}\bigr), \\
  &p_{i,j}^{(m+1)} = p_{i,j}^{(m)}
    + \omega\bigl(\tilde{p}_{i,j} - p_{i,j}^{(m)}\bigr).
\end{split}
\end{equation}

При $\omega = 1$ метод совпадает с Гаусса–Зейделя. При $1 < \omega < 2$ выполняется верхняя релаксация (over-relaxation), ускоряющая сходимость. Оптимальное значение~$\omega$ зависит от размера сетки и коэффициентов системы; для задач смазочного слоя типичный диапазон составляет $\omega \in [1{,}2;\; 1{,}9]$, и на практике параметр подбирается эмпирически.

\textit{Критерий сходимости.}
Итерационный процесс останавливается при выполнении одного из двух критериев. Относительная поправка:
\begin{equation}\label{eq:convergence_change}
  \delta^{(m)} = \frac{\displaystyle\sum_{i,j} \bigl|p_{i,j}^{(m+1)} - p_{i,j}^{(m)}\bigr|}
       {\displaystyle\sum_{i,j} \bigl|p_{i,j}^{(m+1)}\bigr| + \varepsilon_0} < \varepsilon,
\end{equation}
\begin{conditions}
$\sum_{i,j}$ -- суммирование ведётся по узлам, в которых давление является неизвестным;\\
$\varepsilon_0$ -- малая добавка для предотвращения деления на нуль (например, $\varepsilon_0 = 10^{-8}$).
\end{conditions}
Невязка разностного уравнения (вычисляется по внутренним узлам):
\begin{equation}\label{eq:convergence_residual}
\begin{split}
  &r_{i,j}^{(m)} = a_P\,p_{i,j}^{(m)} - a_E\,p_{i+1,j}^{(m)} - a_W\,p_{i-1,j}^{(m)}
    - a_N\,p_{i,j+1}^{(m)} - a_S\,p_{i,j-1}^{(m)} - b_{i,j}, \\
  &\max_{i,j}\bigl|r_{i,j}^{(m)}\bigr| < \varepsilon_r.
\end{split}
\end{equation}

Первый критерий (относительная поправка) проще в реализации и используется в данной работе; второй (невязка) не зависит от параметра релаксации и может применяться как дополнительная проверка. Для снижения накладных расходов проверка сходимости выполняется не на каждой итерации, а с заданным интервалом (например, каждые 500~итераций).

\textit{Инициализация и начальное приближение.}
Простейший вариант~-- нулевое начальное поле давления: $p^{(0)}_{i,j} = 0$. При параметрических расчётах (серия задач с постепенно меняющимися параметрами) в качестве начального приближения используется решение предыдущего шага («тёплый старт»), что существенно сокращает число итераций. Аналогично для нестационарных задач начальным приближением служит решение предыдущего временного шага.

\textit{Вычислительная сложность.}
Стоимость одной итерации составляет $O(N_1 N_2)$ операций (одно обновление на узел). Число итераций до сходимости растёт с размером сетки; для типичных сеток $500 \times 500$ может потребоваться $10^3$-$10^4$~итераций.

Последовательный метод Гаусса–Зейделя (и классический SOR) не допускают параллелизма: обновление узла зависит от уже обновлённых соседей. Параметры итерационного процесса, используемые в данной работе, сведены в таблице~\ref{tab:solver_params}. Для преодоления ограничения на параллелизм применяется шахматная (Red-Black) раскраска сетки, позволяющая обновлять половину узлов одновременно; GPU-реализация этого подхода описана в разделе~3.4.

\begin{table}[!ht]
\centering
\caption{Параметры итерационного процесса}
\label{tab:solver_params}
\begin{tabularx}{\textwidth}{|X|c|}
\hline
\multicolumn{1}{|c|}{\textbf{Параметр}} & \textbf{Значение / диапазон} \\
\hline
Метод               & SOR (Гаусс–Зейдель при $\omega = 1$) \\
Релаксация $\omega$ & $1{,}0$-$1{,}9$ (подбирается эмпирически) \\
Критерий остановки  & относительная поправка $\delta^{(m)} < \varepsilon$ \\
Точность $\varepsilon$ & $10^{-5}$-$10^{-6}$ \\
Интервал проверки   & каждые $10^2$-$10^3$ итераций (типично 500) \\
Макс.\ число итераций & $\leq 50\,000$ (верхняя граница) \\
Инициализация       & $p^{(0)} = 0$ или «тёплый старт» \\
\hline
\end{tabularx}
\end{table}

\FloatBarrier

\textit{Влияние параметра релаксации на сходимость.}
Скорость сходимости SOR критически зависит от выбора
параметра~$\omega$. Теоретически оптимальное значение
определяется спектральным радиусом матрицы итерации Якоби:
$\omega_{\mathrm{opt}} = 2/(1 + \sqrt{1 - \rho^2})$,
где $\rho$~-- спектральный радиус. Для практических задач
смазочного слоя спектральный радиус зависит от размера сетки,
коэффициентов анизотропии ($\Lambda$) и характера зазора;
аналитическое выражение для $\rho$ недоступно. На практике
$\omega$ подбирается следующим образом: для новой конфигурации
проводится несколько пробных расчётов на грубой сетке
с различными $\omega \in [1{,}2;\; 1{,}9]$ и выбирается
значение, минимизирующее число итераций; на мелкой сетке
оптимум, как правило, смещается вверх. Для задач
с микротекстурой, где зазор имеет резкие перепады, типичный
рабочий диапазон $\omega \in [1{,}3;\; 1{,}6]$.
При $\omega \geq 2$ итерационный процесс расходится.

\textit{Роль тёплого старта в параметрических расчётах.}
Значительная часть расчётов в главах 5--8 представляет собой
параметрические серии: $\varepsilon$ меняется с шагом $0{,}1$,
или параметр клина $K$ обходит набор значений. В таких сериях
поле давления $P^{(k-1)}$, полученное при предыдущем значении
параметра, является хорошим начальным приближением для
следующего шага. Тёплый старт снижает число итераций до
сходимости в 3--7 раз по сравнению с нулевым начальным полем,
особенно при малых шагах по параметру. Для расчёта
коэффициентов жёсткости и демпфирования (глава~8), где
требуется многократное решение задачи с близкими геометриями
(возмущение рабочей точки на $\pm\Delta x$ или $\pm\Delta v$),
использование тёплого старта от базового решения снижает
вычислительные затраты аналогично~\cite{sfyris2012,chasalevris2013}.
